% part: intuitionistic-logic
% chapter: semantics
% section: topological-semantics

\documentclass[../../../include/open-logic-section]{subfiles}

\begin{document}

\olfileid{int}{sem}{top}

\olsection{拓扑语义}

为直觉主义逻辑提供语义的另一种方法是使用拓扑这一数学概念。

\begin{defn}
  设 $X$ 是一个集合。$X$ 上的一个\emph{拓扑}是指满足以下性质的集合 $\Top{O}
  \subseteq \Pow{X}$。
  $\Top{O}$ 的元素称为该拓扑的\emph{开集}。集合 $X$ 连同 $\Top{O}$ 一起称为一个
  \emph{拓扑空间}。
  \begin{enumerate}
  \item 空集和整个空间是开集：$\emptyset$, $X \in
    \Top{O}$。
  \item 开集对有限交封闭：如果 $U$, $V \in
    \Top{O}$，那么 $U \cap V \in \Top{O}$。
  \item 开集对任意并封闭：如果对所有 $i \in I$ 都有 $U_i \in
    \Top{O}$，那么 $\bigcup \Setabs{U_i}{i \in
      I} \in \Top{O}$。
  \end{enumerate}
\end{defn}

如果从上下文可以推断出开集族，我们有时也直接用 $X$ 来表示一个拓扑；但需注意，只有当 $X$ 被赋予开集之后，它才能被称为一个拓扑。

\begin{defn}
  直觉主义命题逻辑的一个\emph{拓扑模型}是一个三元组 $\mModel{X} = \tuple{X, \Top{O}, V}$，其中 $\Top{O}$ 是 $X$ 上的一个拓扑，而 $V$ 是一个函数，为每个命题变元指派 $\Top{O}$ 中的一个开集。

  给定一个拓扑模型~$\mModel{X}$，我们可以归纳地定义 $\Prop{X}{!A}$ 如下：
  \begin{enumerate}
  \item $\Prop{X}{\lfalse} = \emptyset$
  \item $\Prop{X}{p} = V(p)$
  \item $\Prop{X}{!A \land !B} = \Prop{X}{!A} \cap \Prop{X}{!B}$
  \item $\Prop{X}{!A \lor !B} = \Prop{X}{!A} \cup \Prop{X}{!B}$
  \item $\Prop{X}{!A \lif !B} = \Interior{(X \setminus \Prop{X}{!A}) \cup
    \Prop{X}{!B}}$
  \end{enumerate}
  这里，$\Interior{V}$ 是将子集 $V \subseteq X$ 映射到其\emph{内部}的函数，即包含于 $V$ 的所有开集的并。换言之，
  \[
  \Interior{V} = \bigcup \Setabs{U}{U \subseteq V \text{ 且 } U \in \Top{O}}.
  \]
\end{defn}

注意，任何集合的内部总是开集，因为它是开集的并。因此，$\Prop{X}{!A}$ 总是一个开集。

尽管拓扑语义非常抽象，但有一些直观的思考方式可以为其提供动机。假设 $X$ 的元素或“点”是使得语句能够被赋值的点。使得 $!A$ 为真的所有点的集合，就是 $!A$ 所表达的命题。并非每个点集都是一个潜在的命题；只有 $\Top{O}$ 中的元素才是。$!A$ 语义地后承 $!B$ 当且仅当在 $!A$ 为真的每个点 $!B$ 也为真，即对所有~$X$ 有 $\Prop{X}{!A}
\subseteq \Prop{X}{!B}$。荒谬语句~$\lfalse$ 从不为真，因此 $\Prop{X}{\lfalse} = \emptyset$。

为了使直觉主义有效法则成立，即为了使 $!A \Proves !B$ 当且仅当 $\Prop{X}{!A} \subset \Prop{X}{!B}$，$!B \land !C$、$!B \lor !C$ 以及 $!B \lif !C$ 所表达的命题必须与 $!B$、$!C$ 所表达的命题具有怎样的关系？我们要求对任意 $!A$ 有 $\lfalse \Proves !A$，这之所以成立，是因为对所有 $U$ 都有 $\emptyset \subseteq U$。由于 $!B \land !C \Proves !B$，我们要求 $\Prop{X}{!B \land !C} \subseteq \Prop{X}{!B}$，类似地有 $\Prop{X}{!B \land !C} \subseteq \Prop{X}{!C}$。满足 $W \subseteq U$ 且 $W \subseteq V$ 的最大集合是 $U \cap V$。反过来，$!B \Proves !B \lor !C$ 且 $!C \Proves !B \lor !C$，因此我们要求 $\Prop{X}{!B} \subseteq \Prop{X}{!B \lor !C}$ 且 $\Prop{X}{!C} \subseteq \Prop{X}{!B \lor !C}$。满足 $U \subseteq W$ 且 $V \subseteq W$ 的最小集合~$W$ 是 $U \cup V$。

$\lif$ 的定义则比较棘手：$!A \lif !B$ 表达的是与 $!A$ 结合后能语义地后承 $!B$ 的最弱的命题。$!A \lif !B$ 与 $!A$ 结合后能语义地后承~$!B$，这一点从 $(!A \lif !B)
\land !A \Proves !B $ 可以清楚看出。因此，$\Prop{X}{!A \lif !B}$ 应该是满足 $\Prop{X}{!A \lif !B} \cap \Prop{X}{!A}
\subset \Prop{X}{!B}$ 的最大的开集，这便引出了我们的定义。

\end{document}